\documentclass{report}
\usepackage{amsmath, amssymb, amsthm}
\usepackage{geometry}
\geometry{a4paper, margin=0.8in}
\usepackage{float}
\usepackage{graphicx}
\usepackage{tikz}
\usetikzlibrary{calc, decorations.pathreplacing}
\usepackage{pgfplots}
\pgfplotsset{compat=1.18}
\usepackage{xcolor}
\usepackage{tcolorbox}
\tcbuselibrary{theorems}
\usepackage{eurosym}
\usepackage{fancyhdr}
\usepackage{booktabs}
\usepackage{listings}

\lstdefinestyle{tcblatex}{
    language=[Visual]Basic,
    basicstyle=\ttfamily\footnotesize,
    keywordstyle=\bfseries,
    commentstyle=\itshape,
    stringstyle=\ttfamily,
    showstringspaces=false,
    breaklines=true,
    columns=fullflexible
}

\begin{document}

\pagestyle{fancy}
\fancyhf{}
\rhead{Report 6 \;|\; Stochastic Methods for Finance \;|\; \thepage}
\renewcommand{\headrulewidth}{0.4pt}
\setlength{\headheight}{14pt}
\thispagestyle{plain} 

% --- PROFESSIONAL HEADER ---
\begin{flushleft}
    {\large \textbf{Stochastic Methods for Finance - Report 6}} \\
    \rule{\textwidth}{0.4pt} \\ 
    \textbf{Analyst:} Davide Quartucci \hfill \textbf{Date:} April 20, 2026 \\
    \textbf{Underlying Basket:} Deutsche Lufthansa AG \;|\; Air France KLM SA \;|\; Ryanair Holdings PLC \\
\end{flushleft}

\begin{center}
    \fcolorbox{black}{gray!10}{
        \begin{minipage}{0.95\textwidth}
            \small \textit{Computational architecture and numerical routines are available in the supplementary file:} \texttt{Report6\_Quartucci\_Davide.xlsm}
        \end{minipage}
    }
\end{center}

\section{Environment Setup and Baseline Parameters}
The primary objective of this framework is to establish a robust computational environment for stochastic option pricing. The architecture integrates Microsoft Excel's front-end capabilities with a custom-built Visual Basic for Applications (VBA) Monte Carlo engine.

\subsection{Parameter Initialization}
The baseline parameters for the Black-Scholes framework are standardized as follows:
\begin{itemize}
    \item Spot price: $S_0 = 100$
    \item Annualized volatility: $\sigma = 20\%$
    \item Risk-free rate: $r = 1\%$
    \item Time to maturity: $T = 1$ year
    \item Strike price: $K = 99$
\end{itemize}
Time discretization is set to $dt = 1/252$, accurately reflecting the effective trading days within a standard financial year.

\subsection{Stochastic Driver Generation}
Due to the absence of a native normal distribution generator in VBA, the \textbf{Box-Muller} transform was implemented. This algorithm efficiently maps pairs of independent uniformly distributed variables (generated via \texttt{Rnd()}) into standard normal variables $\mathcal{N}(0,1)$, which serve as the fundamental Brownian increments for the simulations.

\section{Computational Architecture and Implementation Methodology}
A critical aspect of this project is the design of the computational infrastructure. The model was engineered to maximize execution speed, maintainability, and clarity, adopting a hybrid approach between spreadsheet functions and object-oriented programming.

\subsection{Excel-VBA Integration Strategy}
The architecture deliberately separates static historical calibration from dynamic stochastic simulation to optimize computational resources:
\begin{itemize}
    \item \textbf{Data Processing (Front-End):} The extraction of historical prices, the computation of daily logarithmic returns, and the estimation of annualized volatilities and correlation matrices ($\rho$) were performed natively in Excel. Utilizing built-in array functions (e.g., \texttt{STDEV.S}, \texttt{CORREL}) provides immediate inspectability of the market inputs without running external scripts.
    \item \textbf{Stochastic Engine (Back-End):} The Monte Carlo simulations, which require heavy iterative loops and memory allocation, were exclusively coded in VBA. The macros read the calibrated inputs directly from specific cell coordinates, ensuring a seamless data pipeline.
\end{itemize}

\subsection{Algorithmic Optimization and Modularity}
To achieve institutional-grade performance during the simulation of $N=10{,}000$ paths, several optimization techniques were implemented within the VBA environment:
\begin{enumerate}
    \item \textbf{Pre-computation of Constants:} Variables that remain constant across the time steps—such as the deterministic drift $(r - 0.5\sigma^2)dt$ and the diffusion scalar $\sigma\sqrt{dt}$—were strictly calculated outside the main Monte Carlo loops.
    \item \textbf{Code Modularity:} The pricing engine was structured into independent, self-contained subroutines (Macros) for each derivative class (Vanilla, Asian, Lookback, Worst-Of). This modular design prevents variable cross-contamination and allows discrete execution and debugging for each specific payoff structure.
\end{enumerate}

\subsection{Cholesky Factorization Design}
For the multi-asset framework, the Cholesky decomposition of the $3 \times 3$ correlation matrix was explicitly hardcoded using fundamental algebraic equations rather than employing a generic $n \times n$ iterative loop. This choice drastically reduces the algorithmic complexity and execution time for low-dimensional matrices, guaranteeing the instantaneous generation of the lower triangular matrix $L$ required for the correlated shocks.


\section{Geometric Brownian Motion (GBM) Dynamics}
This section evaluates the engine's capability to project the underlying asset's trajectory over time.

\subsection{Model Specification}
The asset price dynamics follow a Geometric Brownian Motion. To guarantee strictly positive sample paths and eliminate discretization bias, the closed-form exact solution of the Stochastic Differential Equation (SDE) is utilized:
\begin{equation}
    S_{t+dt} = S_t \exp\left[\left(r - \frac{1}{2}\sigma^2\right)dt + \sigma \sqrt{dt} Z\right]
\end{equation}

\subsection{Simulation Output}
Executing $N=100$ independent paths over 252 discrete steps yields a distribution spectrum that materializes market uncertainty.

\begin{figure}[htbp]
    \centering
    \includegraphics[width=0.8\textwidth]{Picture1.png}
    \caption{Monte Carlo projection of 100 GBM paths ($S_0=100, \sigma=20\%, r=1\%, T=1$).}
    \label{fig:gbm}
\end{figure}

The graphical output (Figure \ref{fig:gbm}) confirms that the terminal price dispersion is structurally consistent with the $20\%$ volatility input, bounding the majority of terminal states within the $[60, 150]$ interval.

\section{Vanilla Derivatives Valuation}
This section validates the numerical accuracy of the Monte Carlo engine by pricing European Call and Put options against the Black-Scholes closed-form benchmark.

\subsection{Numerical Schemes}
Three distinct computational approaches were implemented:
\begin{enumerate}
    \item \textbf{1-Step Method:} Leverages the path-independence of European payoffs by executing a single stochastic jump to maturity $T$.
    \item \textbf{Euler-Maruyama Multi-Step:} Discretizes the path into 252 steps via $S_{t+dt} = S_t(1 + r dt + \sigma \sqrt{dt} Z)$. It introduces a known truncation error of order $\mathcal{O}(dt)$.
    \item \textbf{Exact Multi-Step:} Iteratively applies the exact GBM log-normal solution across 252 steps, nullifying local discretization bias.
\end{enumerate}

\subsection{Empirical Convergence}
Based on $N=10{,}000$ iterations, the pricing outcomes demonstrate robust convergence toward the analytical benchmark.

\begin{table}[htbp]
    \centering
    \begin{tabular}{lcccccc}
        \toprule
        & \multicolumn{3}{c}{\textbf{Call}} & \multicolumn{3}{c}{\textbf{Put}} \\
        \cmidrule(lr){2-4} \cmidrule(lr){5-7}
        \textbf{Method} & \textbf{MC Price} & \textbf{SE} & \textbf{Abs Error} & \textbf{MC Price} & \textbf{SE} & \textbf{Abs Error} \\
        \midrule
        1-Step (Closed Form) & 8.943 & 0.137 & 0.024 & 6.906 & 0.096 & 0.026 \\
        Euler Multi-Step & 9.066 & 0.139 & 0.147 & 6.951 & 0.097 & 0.018 \\
        Exact Multi-Step & 8.763 & 0.137 & 0.155 & 7.052 & 0.097 & 0.119 \\
        \midrule
        \textbf{Black \& Scholes} & \textbf{8.918} & - & - & \textbf{6.933} & - & - \\
        \bottomrule
    \end{tabular}
    \caption{Comparative pricing summary for European Vanilla Options ($N=10{,}000$).}
    \label{tab:vanilla_comparison}
\end{table}

The analytical evidence suggests that for path-independent payoffs, the 1-Step approach is computationally dominant: it achieves benchmark-level accuracy while drastically minimizing processing overhead. Multi-step schemes remain mandatory solely for path-dependent structures.

\section{Path-Dependent Structures: Asian and Lookback Options}
Expanding the framework's complexity, this section evaluates exotic derivatives characterized by strong path-dependency.

\subsection{Asian Options (Arithmetic Mean)}
The payoff is determined by the arithmetic average of the underlying's daily closing prices. 

\begin{table}[htbp]
    \centering
    \begin{tabular}{lcccc}
        \toprule
        \textbf{Derivative} & \textbf{MC Price} & \textbf{MC Std Deviation} & \textbf{MC Std Error} \\
        \midrule
        Asian Call & 5.348 & 7.836 & 0.078 \\
        Asian Put & 3.860 & 5.701 & 0.057 \\
        \bottomrule
    \end{tabular}
    \caption{Pricing output for Arithmetic Asian Options ($N=10{,}000$).}
    \label{tab:asian}
\end{table}

The pricing metrics (Table \ref{tab:asian}) quantify the \textit{variance reduction} phenomenon inherent in Asian options. The averaging mechanism dampens volatility spikes, structurally depressing the derivative's fair value compared to its Vanilla equivalent.

\subsection{Lookback Options (Floating Strike)}
Lookback options grant the holder the right to exercise at the most favorable price achieved during the life of the option. The algorithmic implementation tracks the running extrema ($\max(S_t)$ and $\min(S_t)$) at each discrete step.

\begin{table}[htbp]
    \centering
    \begin{tabular}{lcccc}
        \toprule
        \textbf{Derivative} & \textbf{MC Price} & \textbf{MC Std Deviation} & \textbf{MC Std Error} \\
        \midrule
        Lookback Call (Max - $S_T$) & 15.429 & 10.672 & 0.106 \\
        Lookback Put ($S_T$ - Min) & 14.977 & 13.805 & 0.138 \\
        \bottomrule
    \end{tabular}
    \caption{Pricing output for Floating Strike Lookback Options ($N=10{,}000$).}
    \label{tab:lookback}
\end{table}

As illustrated in Table \ref{tab:lookback}, the embedded \textit{hindsight} optionality carries a substantial premium, making Lookback structures significantly more expensive than standard derivatives.

\section{Multi-Asset Pricing: The Worst-Of Certificate}
The final model transitions from a unidimensional SDE to a multidimensional framework to price a structured \textit{Worst-Of} product on an equity basket.

\subsection{Market Calibration and Correlation Matrix}
A sector-specific basket comprising three European airlines (Lufthansa, Air France, Ryanair) was analyzed. Based on a 252-day historical rolling window, the annualized volatilities and the empirical correlation matrix of logarithmic returns were extracted.

\begin{table}[htbp]
    \centering
    \begin{tabular}{lccc}
        \toprule
        \textbf{Asset} & \textbf{LHAG.DE} & \textbf{AIRF.PA} & \textbf{RYA.I} \\
        \midrule
        \textbf{LHAG.DE (Vol: 37.25\%)} & 1.0000 & 0.4565 & 0.4791 \\
        \textbf{AIRF.PA (Vol: 50.62\%)} & 0.4565 & 1.0000 & 0.5279 \\
        \textbf{RYA.I (Vol: 30.31\%)}   & 0.4791 & 0.5279 & 1.0000 \\
        \bottomrule
    \end{tabular}
    \caption{Empirical Correlation Matrix ($\rho$) and Annualized Volatilities.}
    \label{tab:correlation}
\end{table}

\subsection{Dependence Modeling via Cholesky Decomposition}
To bypass the assumption of independence, the correlation matrix $\rho$ was factorized using the Cholesky decomposition algorithm, yielding a lower triangular matrix $L$. At each increment, multiplying $L$ by a vector of standard independent normals generates the cross-sectionally correlated stochastic shocks driving the multidimensional simulation.

\subsection{Derivative Valuation and Risk Interpretation}
The priced instrument is a \textit{Worst-Of Bonus Certificate} featuring a European barrier at 60\% of the initial spot prices and a 5\% fixed premium conditional on barrier integrity. 

\begin{table}[htbp]
    \centering
    \begin{tabular}{lc}
        \toprule
        \textbf{Parameter} & \textbf{Value} \\
        \midrule
        Nominal Value & \euro{100} \\
        Barrier Level & 60\% \\
        \midrule
        \textbf{Fair Value (Price)} & \textbf{\euro{87.00}} \\
        MC Std Deviation & 26.56 \\
        MC Std Error & 0.265 \\
        \bottomrule
    \end{tabular}
    \caption{Valuation metrics for the Worst-Of Bonus Certificate ($N=10{,}000$).}
    \label{tab:worstof}
\end{table}

The resulting Fair Value of \euro{87.00} (Table \ref{tab:worstof}) rigorously prices the combined risk factors. The extreme historical volatility of the underlying assets (e.g., Air France at $>50\%$), coupled with the severe multiplicative risk inherent in the "Worst-Of" payout structure, generates a high statistical probability of a \textit{Knock-In} event. The positive correlation ($\approx 0.50$) provides only marginal mitigation, ultimately justifying the significant discount against the nominal value.

\section{Conclusion and Executive Summary}
The development and execution of this quantitative framework successfully demonstrated the flexibility and numerical stability of Monte Carlo methods in pricing diverse financial derivatives. 

The empirical results confirm several key theoretical principles:
\begin{itemize}
    \item \textbf{Computational Efficiency:} For path-independent structures (European Vanillas), the 1-Step closed-form solution is computationally superior, avoiding unnecessary loops without sacrificing accuracy. Euler and Exact multi-step schemes become strictly necessary only when dealing with path-dependent payoffs.
    \item \textbf{Exotic Premia:} Path-dependency fundamentally alters the risk profile of a derivative. Averaging mechanisms (Asian options) depress the option's value via variance reduction, whereas extreme-value tracking (Lookback options) commands a significant premium due to hindsight optionality.
    \item \textbf{Multidimensional Risk:} Transitioning from single-asset to multi-asset pricing requires robust correlation modeling. The implementation of the Cholesky factorization effectively captured the dependence structure of the aviation basket. The valuation of the Worst-Of Certificate highlighted how extreme idiosyncratic volatility and structural "Worst-Of" conditions drastically discount the derivative's fair value, pricing in a high probability of barrier breach.
\end{itemize}
Overall, the VBA-Excel integrated engine proved to be a reliable and highly optimized infrastructure for stochastic derivative pricing.
\end{document}
